%% 10_appendixes_00.tex

\chapter{\acrshort{dsrm} Delivered Artifacts}
\label{chap:artifacts-list}

\section{List of Delivered Artifacts}
\label{appex:dsrm-artifacts}

% Below is a structured overview of the primary artifacts associated with each \acrshort{dsrm} phase.

\noindent\rule{13cm}{0.3pt}

\paragraph{Artifact: Master Thesis Proposal}
\begin{description}
  \item[\textbf{Status:}] Completed
  \item[\textbf{Format:}] Document
  \item[\textbf{Phase:}] Problem Identification and Motivation, Define Objectives for a Solution
  \item[\textbf{Activities:}] Initial meetings with stakeholders, shadowing existing data workflows and methods
  \item[\textbf{Methods:}] Software Requirements Engineering (main), \acrshort{slr} (main)
  \item[\textbf{Acceptance Criteria:}] Approval from external stakeholders and supervisors, master thesis proposal acceptance
  \item[\textbf{Stakeholders:}] Supervisors
  \item[\textbf{Tentative Dates:}] 24.06.2024 – 24.12.2024
  \item[\textbf{Estimated Effort:}] 5 months
\end{description}

\noindent\rule{13cm}{0.3pt}

\paragraph{Artifact: Preliminary Table of Content}
\begin{description}
  \item[\textbf{Status:}] Completed
    \item[\textbf{Format:}] Document
  \item[\textbf{Phase:}] Problem Identification and Motivation, Define Objectives for a Solution
  \item[\textbf{Activities:}] Initial meetings with stakeholders, shadowing existing data workflows and methods, and \acrshort{slr}.
  \item[\textbf{Methods:}] Software Requirements Engineering (main), \acrshort{slr} (main)
  \item[\textbf{Acceptance Criteria:}] Approval from external stakeholders and supervisors, master thesis proposal acceptance
  \item[\textbf{Stakeholders:}] Supervisors
  \item[\textbf{Tentative Dates:}] 01.01.2025 – 02.02.2025
  \item[\textbf{Estimated Effort:}] 1 month
  \end{description}

\noindent\rule{13cm}{0.3pt}

\paragraph{Artifact: Preliminary Project Plan Overview (Appendix ~\ref{appex:wbs})}
\begin{description}
  \item[\textbf{Status:}] Completed
  \item[\textbf{Format:}] Document
  \item[\textbf{Phase:}] Define Objectives for a Solution, and Design and Development \& Artifact Development
  \item[\textbf{Activities:}] Project plan overview design
  \item[\textbf{Methods:}] \acrshort{scrum} (main), \acrshort{crisp} (main), \acrshort{dsrm} (secondary)
  \item[\textbf{Acceptance Criteria:}] Evidence of addressing research questions, alignment with research deadlines, supervisor approval, research methods seminar paper acceptance
  \item[\textbf{Stakeholders:}] Supervisors, Research Methods Seminar Professor
  \item[\textbf{Tentative Dates:}] 01.01.2025 – 02.02.2025
  \item[\textbf{Estimated Effort:}] 1 month
\end{description}

\noindent\rule{13cm}{0.3pt}

\paragraph{Artifact: Preliminary Literature Review (Appendix ~\ref{chap:slr-report})}
\begin{description}
  \item[\textbf{Status:}] Completed
    \item[\textbf{Format:}] Document
  \item[\textbf{Phase:}] Problem Identification and Motivation, and Define Objectives for a Solution
  \item[\textbf{Activities:}] Initial meetings with stakeholders, shadowing existing data workflows and methods, and \acrshort{slr}.
  \item[\textbf{Methods:}] \acrshort{slr} (main)
  \item[\textbf{Acceptance Criteria:}] Approval from external supervisors
  % \item[\textbf{Stakeholders:}] External and internal supervisors
  \item[\textbf{Stakeholders:}] Supervisors
  \item[\textbf{Tentative Dates:}] 01.01.2025 – 02.02.2025
  \item[\textbf{Estimated Effort:}] 1 month
\end{description}

\noindent\rule{13cm}{0.3pt}

\paragraph{Artifact: Prototype Development}
\begin{description}
  \item[\textbf{Status:}] Completed
  \item[\textbf{Format:}] Working Software
  \item[\textbf{Phase:}] Define Objectives for a Solution, Design and Development
  \item[\textbf{Activities:}] Software development
  \item[\textbf{Methods:}] \acrshort{scrum} (main), \acrshort{slr} (main), \acrshort{crisp} (secondary)
  \item[\textbf{Acceptance Criteria:}] Supervisor approval, preprocessing implemented, bioacoustic methods implemented, initial model training and testing
  % \item[\textbf{Stakeholders:}] EPFL Supervisor
  \item[\textbf{Stakeholders:}] Supervisor
  \item[\textbf{Tentative Dates:}] 01.10.2024 – 08.01.2025
  \item[\textbf{Estimated Effort:}] 3 months
\end{description}

\noindent\rule{13cm}{0.3pt}

\paragraph{Artifact: FOREST - Application Development}
\begin{description}
  \item[\textbf{Status:}] In-Progress
  \item[\textbf{Format:}] Working Software
  \item[\textbf{Phase:}] Demonstration and Implementation
  \item[\textbf{Activities:}] \acrshort{scrum} with bi-weekly sprints, modular implementation
  \item[\textbf{Methods:}] \acrshort{scrum} (main), \acrshort{crisp} (main), Clean Architecture (secondary), \acrshort{slr} (secondary)
  \item[\textbf{Acceptance Criteria:}] Supervisor approval, evidence of addressing research questions
  % \item[\textbf{Stakeholders:}] EPFL \& TU Wien Supervisors
  \item[\textbf{Stakeholders:}] Supervisors
  \item[\textbf{Tentative Dates:}] 11.11.2024 – 09.09.2025
  \item[\textbf{Estimated Effort:}] 10 months
\end{description}

\noindent\rule{13cm}{0.3pt}

\paragraph{Artifact: FOREST - Preprocessing Tasks for Time Series Data}
\begin{description}
  \item[\textbf{Status:}] Completed
  \item[\textbf{Format:}] Dataset
  \item[\textbf{Phase:}] Demonstration and Implementation
  \item[\textbf{Activities:}] \acrshort{scrum} with bi-weekly sprints
  \item[\textbf{Methods:}] \acrshort{scrum} (main), \acrshort{crisp} (main), Clean Architecture (secondary), \acrshort{slr} (secondary)
  \item[\textbf{Acceptance Criteria:}] Supervisor approval, evidence of addressing research questions
  % \item[\textbf{Stakeholders:}] EPFL Supervisors
  \item[\textbf{Stakeholders:}] Supervisors
  \item[\textbf{Tentative Dates:}] 01.11.2024 – 01.12.2024
  \item[\textbf{Estimated Effort:}] 2 months
\end{description}

\noindent\rule{13cm}{0.3pt}

\paragraph{Artifact: FOREST - Apply Bioacoustic Metrics}
\begin{description}
  \item[\textbf{Status:}] Completed
  \item[\textbf{Format:}] Dataset
  \item[\textbf{Phase:}] Demonstration and Implementation
  \item[\textbf{Activities:}] \acrshort{scrum} with bi-weekly sprints
  \item[\textbf{Methods:}] \acrshort{scrum} (main), \acrshort{crisp} (main), \acrshort{slr} (main)
  \item[\textbf{Acceptance Criteria:}] Supervisor approval, evidence of addressing research questions
  % \item[\textbf{Stakeholders:}] EPFL Supervisors
  \item[\textbf{Stakeholders:}] Supervisors
  \item[\textbf{Tentative Dates:}] 11.11.2024 – 10.01.2025
  \item[\textbf{Estimated Effort:}] 2 months
\end{description}

\noindent\rule{13cm}{0.3pt}

\paragraph{Artifact: FOREST - Apply MIR Techniques for Data Transformation}
\begin{description}
  \item[\textbf{Status:}] Discarded (Out-of-Scope)
  \item[\textbf{Format:}] Dataset
  \item[\textbf{Phase:}] Demonstration and Implementation
  \item[\textbf{Activities:}] \acrshort{scrum} with bi-weekly sprints
  \item[\textbf{Methods:}] \acrshort{scrum} (main), \acrshort{crisp} (main), \acrshort{slr} (main)
  \item[\textbf{Acceptance Criteria:}] Supervisor approval, evidence of addressing research questions
  % \item[\textbf{Stakeholders:}] EPFL Supervisors
  \item[\textbf{Stakeholders:}] Supervisors
  \item[\textbf{Tentative Dates:}] - % 04.04.2025 – 05.05.2025
  \item[\textbf{Estimated Effort:}] - % 1 month
\end{description}

\noindent\rule{13cm}{0.3pt}

\paragraph{Artifact: FOREST - Curated Dataset}
\begin{description}
  \item[\textbf{Status:}] Completed
  \item[\textbf{Format:}] Dataset
  \item[\textbf{Phase:}] Demonstration and Implementation
  \item[\textbf{Activities:}] \acrshort{scrum} with bi-weekly sprints
  \item[\textbf{Methods:}] \acrshort{crisp} (main), \acrshort{slr} (main)
  \item[\textbf{Acceptance Criteria:}] Supervisor approval, able to be used to train the models, evidence of addressing research questions
  % \item[\textbf{Stakeholders:}] EPFL Supervisors
  \item[\textbf{Stakeholders:}] Supervisors
  \item[\textbf{Tentative Dates:}] 11.11.2024 – 03.03.2025
  \item[\textbf{Estimated Effort:}] 3 months
\end{description}



\noindent\rule{13cm}{0.3pt}

\paragraph{Artifact: Model Architecture}
\begin{description}
  \item[\textbf{Status:}] Completed
  \item[\textbf{Format:}] PyTorch Objects
  \item[\textbf{Phase:}] Demonstration and Implementation, and Evaluation
  \item[\textbf{Activities:}] \acrshort{scrum} with bi-weekly sprints
  \item[\textbf{Methods:}] \acrshort{scrum} (main), \acrshort{crisp} (main), \acrshort{slr} (main)
  \item[\textbf{Acceptance Criteria:}] Supervisor approval, Evidence of decent performance time, and prediction scores.
  % \item[\textbf{Stakeholders:}] EPFL Supervisors
  \item[\textbf{Stakeholders:}] Supervisors
  \item[\textbf{Tentative Dates:}] 12.12.2024 – 13.03.2025  
\item[\textbf{Estimated Effort:}] 3 months
\end{description}

\noindent\rule{13cm}{0.3pt}


\paragraph{Artifact: Model Training \& Testing and Evaluation}
\begin{description}
  \item[\textbf{Status:}] Completed
  \item[\textbf{Format:}] Performance Results JSON
  \item[\textbf{Phase:}] Demonstration and Implementation, and Evaluation
  \item[\textbf{Activities:}] \acrshort{scrum} with bi-weekly sprints
  \item[\textbf{Methods:}] \acrshort{crisp} (main)
  \item[\textbf{Acceptance Criteria:}] Supervisor approval, Evidence of decent performance time, and prediction scores.
  % \item[\textbf{Stakeholders:}] EPFL Supervisors
  \item[\textbf{Stakeholders:}] Supervisors
  \item[\textbf{Tentative Dates:}] 12.12.2024 – 13.03.2025
  \item[\textbf{Estimated Effort:}] 3 months
\end{description}

\noindent\rule{13cm}{0.3pt}

\paragraph{Artifact: Models' Performance Evaluation with External Datasets}
\begin{description}
  \item[\textbf{Status:}] Pending (TBC)
  \item[\textbf{Format:}] Performance Results JSON
  \item[\textbf{Phase:}] Evaluation
  \item[\textbf{Activities:}] \acrshort{scrum} with bi-weekly sprints
  \item[\textbf{Methods:}] \acrshort{scrum} (main), \acrshort{crisp} (main), \acrshort{slr} (main)
  \item[\textbf{Acceptance Criteria:}] Supervisor approval, Evidence of decent performance time, and prediction scores.
  % \item[\textbf{Stakeholders:}] EPFL Supervisors
  \item[\textbf{Stakeholders:}] Supervisors
  \item[\textbf{Tentative Dates:}] 05.05.2025 – 08.08.2025
  \item[\textbf{Estimated Effort:}] 3 months
\end{description}

\noindent\rule{13cm}{0.3pt}


\paragraph{Artifact: Models' Performance Comparison with SOTA Models}
\begin{description}
  \item[\textbf{Status:}] Completed
  \item[\textbf{Format:}] Performance Results JSON
  \item[\textbf{Phase:}] Evaluation
  \item[\textbf{Activities:}] \acrshort{scrum} with bi-weekly sprints
  \item[\textbf{Methods:}] \acrshort{scrum} (main), \acrshort{crisp} (main), \acrshort{slr} (main)
  \item[\textbf{Acceptance Criteria:}] Supervisor approval, Evidence of decent performance time, and prediction scores.
  % \item[\textbf{Stakeholders:}] EPFL Supervisors
  \item[\textbf{Stakeholders:}] Supervisors
  \item[\textbf{Tentative Dates:}] 12.02.2025 – 05.05.2025
  \item[\textbf{Estimated Effort:}] 3 months
\end{description}

\noindent\rule{13cm}{0.3pt}


\paragraph{Artifact: Master Thesis Manuscript and Presentations}
\begin{description}
  \item[\textbf{Status:}] In-Progress
  \item[\textbf{Format:}] Document
  \item[\textbf{Phase:}] Communication and Research Contribution
  \item[\textbf{Activities:}] Scientific writing, \acrshort{slr}
  \item[\textbf{Methods:}] Scientific Writing (main), \acrshort{slr} (main), \acrshort{crisp} (secondary)
  \item[\textbf{Acceptance Criteria:}] Supervisors approval, evidence of addressing research questions
  % \item[\textbf{Stakeholders:}] EPFL \& TU Wien Supervisors, TU Wien Informatics Decanature
  \item[\textbf{Stakeholders:}] TU Wien Informatics Decanature
  \item[\textbf{Tentative Dates:}] 11.11.2024 – 09.09.2025
  \item[\textbf{Estimated Effort:}] 10 months
\end{description}

\noindent\rule{13cm}{0.3pt}


% \newpage


\begin{comment}

\subsubsection{\acrfull{dsrm}}

% \todo{include what was actually done during this phase}

This well-known research framework is a problem-solving approach focused on the creation, evaluation, and refinement of innovative artifacts to address complex organizational and technological challenges. Rooted in disciplines like engineering and information systems, it emphasises both practical and theoretical advancement. This dual goal is achieved through rigorous evaluation methods such as experiments, case studies, and simulations \cite{dsrm-Engstroem2020}. The well-defined artifacts serve as tangible solutions to identified problems while contributing to the broader knowledge base \cite{dsrm-info-8dbe4657-b14a-3c79-b3fe-5a48b6386abb}.

\acrshort{dsrm} is an iterative methodology, involving steps such as identifying a relevant problem, defining objectives, developing and demonstrating the artifact, evaluating its performance, and communicating results \cite{dsrm-Peffers01122007}.

The \acrshort{dsrm} can be described as a framework composed by mainly six (6) phases: 

\begin{enumerate}
\item \textbf{Problem identification and motivation:}
The first phase focused on exploring theoretical and well adopted knowledge about Service Design and data science. This can be performed by shadowing and interviewing the relevant stakeholders, or potential users of the solution, and their current way of work. Also, the phase proposes to identifying the environment, culture, organisational mission, and vision, as well as the stakeholders (people, organizations, technologies) to frame the research problem.


\item \textbf{Define Objectives for a Solution:}
The main goal should be inferred rationally from the problem definition and knowledge of what is possible and feasible with the existing resources, or how to obtain and transform external resources that can be acquired.

\item \textbf{Design and Development \& Artifact Development:}
This step includes determining the artifact’s functionality and its architecture and then creating the actual artifact to meet the identified business need.

\item \textbf{Demonstration and Implementation:}
Demonstration may involve its use in experimentation, simulation, case studies, or other appropriate activities. The focus is to use the existing resources to deliver an artifact that can, potentially, solve the problem. This phase enhance the constant feasibility of the solution and iterates in continuous improvement.

\item \textbf{Evaluation:}
This phase questions the quality of the output, and how well the artifact supports a solution of the problem. As well as criticises the implemented methods to propose further iterations or in-depth discussions with stakeholders who can support the continuous improvement and contextualise the added value of the tangible proposed solution.\\
The deliverable artifact must be rigorously evaluated with respect to the utility provided in solving the problem, and tested through workshops, within a specific use cases: Ideal, \& hypothetical scenarios, and against predefined performance assessment metrics.

\item \textbf{Communication and Research Contribution}:
The results must be communicated effectively to both technology-oriented and management-oriented audiences.

\end{enumerate}

This is the primary research framework for this master thesis, guiding the development and validation of the deliverable artifacts. Each phase of the \acrshort{dsrm} results in one or more artifacts. The initial phases involve analytical tasks, such as identifying the problem and defining objectives for a feasible solution. These artifacts have been planned in advance to forecast the necessary activities for a successful project execution, ensuring quality, efficiency, and effectiveness. In contrast, the later phases emphasize actionable tasks. Given the software-centric nature of this research, the Design \& Development, and Demonstration \& Implementation phases follow an iterative approach for continuous refinement (see Appendix \ref{chap:artifacts-list}).

\newpage

Each artifact is described using the following template:

\noindent\rule{13cm}{0.3pt}

\paragraph{Artifact Template}
\label{sec:intro-art-template}

\begin{description}[leftmargin=2cm, labelindent=1cm]
  \item[\textbf{Phase:}] The phase within the Design Science Research Methodology, which includes Problem Identification and Motivation, Define Objectives for a Solution, Design and Development (Artifact Development), Demonstration and Implementation, Evaluation, and Communication and Research Contribution.
  
  \item[\textbf{Artifact:}] The tangible deliverable generated as an outcome of the activities.

    \item[\textbf{Format:}] The type of deliverable resulting after executing the activities, such as a document, dataset, or model data. It defines how information is organized, stored, and interpreted to use in further steps of the methodology, or as a final result.
    
  \item[\textbf{Status:}] The current state of the artifact, which can be one of the following: Completed, In-Progress, Pending, To-Be-Confirmed.

  
  \item[\textbf{Activities:}] The specific actions or set of tasks required to produce the artifact.
  
  \item[\textbf{Method:}] The scientific method or framework used to perform the activities.
  
  \item[\textbf{Acceptance Criteria:}] The indicators used to evaluate and approve the deliverable.
  
  \item[\textbf{Stakeholders:}] The individuals or groups involved in executing the activities.
  
  \item[\textbf{Tentative Dates:}] The planned timeline for carrying out the activities.
  
  \item[\textbf{Estimated Effort (optional):}] The expected effort required to complete the activities.
\end{description}

\noindent\rule{13cm}{0.3pt}


\subsubsection{\acrfull{slr}}

% \todo{include what was actually done during this phase}

A \acrshort{slr} as a research method is a structured to collect, assess, and synthesise external research sources. It follows a well-defined step-by-step process that ensures reproducibility, transparency and academic rigour \cite{slr-SNYDER2019333, slr-SCHROER2021526}. \acrshort{slr} are commonly used to answer specific research questions by analysing existing studies, summarise evidence, identify research gaps, and guide further research steps. It has been originated in the medical research. Although, it has been well-adopted among a variety of other domains such as Management Science, Psychology, and Information Systems, among many others \cite{slr-Stupariu2022, slr-DIVAIO2020283, slr-kitchenham2007guidelines}.

Literature reviews are essential for research, serving as a baseline for knowledge development, guiding policy and practice, and inspiring new ideas. They support future research and theory building. However, conducting and evaluating literature reviews can be complex \cite{slr-SNYDER2019333}. Following clear guidelines ensures more rigorous reviews, helping identify genuine research gaps, refine hypotheses, and improve the overall quality of research within the specific topic's domain \cite{slr-SCHROER2021526}.

In particular for this master thesis, the \acrshort{slr} acts as a supportive method, enabling the identification of novel solutions and insights from the academic and scientific community while identifying patterns from similar use case in domains such as Data Science, Information Systems, Biodiversity Conservation, and Bioacoustics (see Appendix \ref{chap:slr-report}).

\subsubsection{Search Strategy}

The Search Strategy process is based on the search of "Topics" among certain domains or "Groups". The combination of these Topics, and Groups will serve to the Articles Extraction Process. After some of the most relevant Articles for the RQs are selected, the further steps will assure that only the most influential Articles will be used for the Master Thesis Research.

\paragraph{STEP 0 - Definition of the Topics and Groups}
\newline

\begin{description}[leftmargin=2cm, labelindent=1cm]

\item[\textbf{Topics:}] Based on the \acrshort{rqs}.

\begin{itemize}
    
    \item [\textbf{RQ1:}] Data Exploratory Analysis, Cluster Analysis, Feature Extraction, Factor Analysis, PCA, ICA, Data Visualisation\\
    
    \item [\textbf{RQ2:}] Preprocessing, Data Transformation, Data Augmentation, Music Information Retrieval, Bioacoustic Indexes\\
    
    \item [\textbf{RQ3:}] Machine Learning Framework, Modelling, Model Architecture, Classification Model, Predictive Model, Deep Learning Model, CNN, RNN, LSTM, Model Evaluation, Performance Evaluation, Classification, Datasets
    
\end{itemize}

\item[\textbf{Groups:}] Domain Specific.

\begin{itemize}
    \item [\textbf{RQ1-3:}] Time Series, Audio Signals, Audio Records, Acoustics, Bioacoustics, Soundscapes, Environmental Audio
\end{itemize}
    
\end{description}

\paragraph{STEP 1 - Extraction of Papers from Library Databases}

\begin{description}[leftmargin=2cm, labelindent=1cm]
    
    \item[\textbf{Data Bases:}] The specific bases that are relevant to the Research questions. The most relevant Databases that will be used for this research project are: Österreichische Dissertationsdatenbank, IEEE/IET Electronic Library (IEL), ACM Digital Library, arXiv, ScienceDirect, Springer Nature Link, and Scopus. 
  
    \item[\textbf{Document Type:}]  The type of document to query. For the most of the search queries they will be either Scientific \textbf{Articles and/or Books.}
  
    \item[\textbf{Timespan:}] The search will be filtered by only the items that are not older than 10 years (2015 - 2025). This is due to the dynamist of the research field. The goal is to retrieve the most relevant advances on the field from current research outcomes.
    
    \item[\textbf{Language:}] English.
  
    \item[\textbf{Query:}] By using the \textbf{Topics} in truncated combinations with the \textbf{Groups} as search strings.
    
\end{description}

\paragraph{STEP 2 - Identification of Relevant Article}

\begin{description}[leftmargin=2cm, labelindent=1cm]
        
    \item[\textbf{Method:}] Scientific reading of the Abstract and the full article.
  
    \item[\textbf{Acceptance Criteria:}] There will be three (3) main requirements.
    % \newline
    \begin{itemize}
        \item The Article addresses, at least partially, one of the RQ.
        \item The Article utilises Time Series Data, or Audio Records.
        \item Involvement of at least four (4) Authors.
    \end{itemize}
    
\end{description}

\paragraph{STEP 3 - Manually Tracking Citations For Additional Relevant Articles}

\begin{description}[leftmargin=2cm, labelindent=1cm]
        
    \item[\textbf{Method:}] Internet Information Search for citation tracking.
  
    \item[\textbf{Acceptance Criteria:}] There will be two (2) main requirements.
    % \newline
    \begin{itemize}
        \item Same Criteria than for STEP 2.
        \item Highly ranked journals.
    \end{itemize}
    
\end{description}


\paragraph{STEP 4 - Identification of Most Influential Articles}

\begin{description}[leftmargin=2cm, labelindent=1cm]
        
    \item[\textbf{Method:}] Scientific reading of the full article.
  
    \item[\textbf{Acceptance Criteria:}] There will be Three (3) main requirements.
    % \newline
    \begin{itemize}
        \item The Article has been recently published.
        \item The Article addresses more than one of the RQs.
        \item The number of citations per year.
    \end{itemize}
    
\end{description}

\noindent\rule{13cm}{0.3pt}
\newpage

\subsubsection{\acrfull{crisp}}

% \todo{include what was actually done during this phase}

The data processing phase will adopt the \acrshort{crisp} (Cross-Industry Standard Process for Data Mining) methodology \cite{crisp-Wirth2000CRISPDMTA, slr-SCHROER2021526, slr-kitchenham2007guidelines}, a widely recognized approach for handling data science projects, and will be tailored to bioacoustic signal processing.

The \acrshort{crisp} is a well-structured framework designed to guide the execution of projects that involve data as source of information \cite{crisp-Chapman2000CRISPDM1S, crisp-Wirth2000CRISPDMTA}. \acrshort{crisp} offers a comprehensive, flexible, and replicable approach for data mining projects. By dividing the process into six (6) clear phases: 

% Business Understanding, Data Understanding, Data Preparation, Modelling, Evaluation, and Deployment \cite{crisp-Chapman2000CRISPDM1S, crisp-Wirth2000CRISPDMTA}.

\begin{enumerate}

\item\textbf{Business Understanding}

This phase focuses on identifying the objectives of the project from a business perspective and converting them into clear data mining goals while alignment of the technical efforts with business priorities:

\begin{itemize}
    \item Determining the main objectives, such as increasing sales or reducing churn.
    \item Assessing the current business situation, including resources, constraints, and risks.
    \item Defining success criteria for both business and data mining goals.
    \item Developing a high-level project plan to align efforts with business needs.
\end{itemize}

\item \textbf{Data Understanding}

This phase involves collecting, exploring, and analysing data to assess its quality and relevance to ensure that the project works with reliable and meaningful data.

\begin{itemize}
    \item Gathering initial datasets and reviewing their characteristics.
    \item Assessing data quality and identifying missing or inconsistent values.
    \item Exploring data distributions to uncover patterns or insights.
    \item Formulating hypotheses based on observed data trends.
\end{itemize}

\item \textbf{Data Preparation}

Raw data is cleaned and transformed into a structured format suitable for modelling. This is a critical phase that impact significantly to achieving reliable results in further phases.

\begin{itemize}
    \item Selecting relevant tables, records, and attributes.
    \item Cleaning the data to address errors and inconsistencies.
    \item Creating new features or variables to enhance model accuracy.
    \item Formatting and integrating data into a unified structure.
\end{itemize}

\item \textbf{Modelling}

The goal of this phase is to select and leverage appropriate algorithms in conjunction with the data to implement a predictive or descriptive models. This phase translates prepared data into actionable insights.

\begin{itemize}
    \item Choosing modelling techniques suited to the problem type.
    \item Setting up test designs to validate model performance.
    \item Building and fine-tuning models.
    \item Assessing model outputs and refining as needed.
\end{itemize}

\item \textbf{Evaluation}

In this phase, models are evaluated to ensure they meet the defined business objectives and can be utilised for practical applications.

\begin{itemize}
    \item Comparing model results against success criteria defined earlier.
    \item Reviewing each step of the process to identify potential gaps.
    \item Assessing the model’s business impact.
    \item Deciding whether to proceed with deployment or revisit earlier phases.
\end{itemize}

\item \textbf{Deployment}

The final phase operationalizes the models and insights for real-world use.

\begin{itemize}
    \item Generating reports or visualizations for stakeholders.
    \item Automating the model for continuous application.
    \item Establishing monitoring and maintenance procedures.
\end{itemize}

\end{enumerate}

Therefore, \acrshort{crisp} provides a roadmap that enhances collaboration, ensures transparency, and increases the likelihood of success. This methodology is widely applicable across industries and can be tailored for specific project needs to ensure tangible deliverables artifacts. In the particular case of this master thesis the deliverable artifacts out of the \acrshort{crisp} framework can be found in the Appendix \ref{chap:crisp-report}.

\subsubsection{Implementation of the Method}

The study follows a structured approach aligned with \acrshort{crisp}:
\begin{enumerate}
    \item \textbf{Business Understanding}
        \begin{itemize}
            \item \textbf{Context:} See Section \ref{sec:motivation}.
            \item \textbf{Research Questions:} See Section \ref{sec:rqs}.
            \item \textbf{Objective:} FOREST will support  the efforts on biodiversity loss monitoring.
            \item \textbf{Requirements, Assumptions, and Constraints:} The research is based on soundscapes recorded at the same time of year and hour of the day to ensure consistency in data collection.
            \item \textbf{Risks and Contingencies:} The timeline for the Master's Thesis Research Project may be extended if necessary, with a maximum duration of 12 months. Currently, an extension to 10 months has already been granted.
            \item \textbf{Costs and Benefits:}  This master's thesis will contribute to the author's degree completion while also supporting a larger research initiative. Additionally, it will aid ongoing efforts to monitor biodiversity loss in the rainforests of Latin America.
            \item \textbf{Data Mining Goals:} Develop an explainable and reliable model capable of predicting the class of soundscapes based on their characteristics.
            \item \textbf{Data Mining Success Criteria:} The model's classification performance, measured using the F1 score after data preprocessing, should be comparable to SOTA models.
            \item \textbf{Produce Project Plan:} The detailed project plan is outlined in Appendix \ref{appex:wbs}.
        \end{itemize}

    \item \textbf{Data Understanding}
        \begin{itemize}
            \item Unsupervised analysis of 60 hours of soundscapes across four regions.
            \item Exploratory Statistical Analysis, Cluster Analysis, PCA, ICA, and Factor Analysis.
        \end{itemize}
    
    \item \textbf{Data Preparation}\\

    Tailored methods for audio signal preprocessing will be applied, drawing from two key domains: bioacoustics and music information retrieval. This phase aims to produce a well-curated dataset that will be used for model training, testing, and evaluating the classification task.
    
        \begin{itemize}
            \item Audio signal processing techniques from bioacoustics and music information retrieval.
            \item Extraction of bioacoustic indexes.
            \item Augmentation methods to enhance dataset quality.
            \item Creation of a curated dataset for model training and evaluation.
        \end{itemize}
        
    \item \textbf{Modelling}\\

    Based on the dataset's characteristics and the \acrshort{slr} process, a combination of model architectures has been explored for the classification task. One model will classify soundscapes based on extracted features, while another will handle the data as a vector or matrix, leveraging the time-series nature of the dataset.

        \begin{itemize}
            \item Design deep learning architectures.
            \item Train and test models.
            \item Evaluate models based on performance metrics.
        \end{itemize}
    
    \item \textbf{Evaluation}\\

    The evaluation will compare the performance of our application (FOREST) in classifying soundscapes using our dataset, as well as external datasets, against SOTA audio classification models. The following tasks will be carried out:
    
        \begin{itemize}
            \item Compare FOREST model against SOTA models with the internally curated dataset.
            \item Conduct empirical experiments.
            \item Benchmark classification performance against SOTA models with external datasets.
        \end{itemize}
        
    \item \textbf{Deployment:}
        \begin{itemize}
            \item Development the software under the Scrum Agile methodology with 2-3 week sprints.
            \item Train and test models on HPC systems (Vienna Scientific Cluster from TU Wien, and SCITAS from EPFL).
            \item Host the project on GitHub, store data on Zenodo, and public access to the FOREST will be deployed under an MIT license at url \href{http://www.soundforest.app/}{www.soundforest.app}.
        \end{itemize}

\end{enumerate}

\end{comment}

% \chapter{\acrfull{slr} Report}
% \label{chap:slr-report}

% \section{Project Plan}
% \label{appex:wbs}